\documentclass[11pt,a4paper]{article}

\usepackage[T1]{fontenc}
\usepackage{lmodern}
\usepackage[margin=2.6cm]{geometry}
\usepackage{amsmath,amssymb,amsthm,mathtools}
\usepackage{microtype}
\emergencystretch=2em
\usepackage{booktabs}
\usepackage[table]{xcolor}
\usepackage{enumitem}
\usepackage{array}
\usepackage{tabularx}
\usepackage[colorlinks=true,linkcolor=blue!50!black,citecolor=blue!50!black,urlcolor=blue!50!black]{hyperref}

\setlist{itemsep=2pt,topsep=4pt}

% Toads and Frogs notation
\newcommand{\Tt}{\mathsf{T}}
\newcommand{\Ff}{\mathsf{F}}
\newcommand{\E}{\square}
\newcommand{\pos}[1]{\ensuremath{#1}}
\newcommand{\lean}[1]{\texttt{#1}}
\newcommand{\new}{\cellcolor{orange!22}}
\newcommand{\fuzzy}{\mathrel{\|}}

\theoremstyle{plain}
\newtheorem{theorem}{Theorem}
\newtheorem{proposition}[theorem]{Proposition}
\newtheorem{lemma}[theorem]{Lemma}
\theoremstyle{remark}
\newtheorem*{remark}{Remark}

\title{New results on Toads and Frogs:\\
Thanatipanonda's Conjectures T1 and T2, the column $b=4$,\\
and counterexamples to his Conjecture T4}
\author{Charlie Reams\\[4pt]
\small Lean development: \url{https://github.com/calr3/toads-and-frogs-erickson}}
\date{November 2026}

\begin{document}
\maketitle

\begin{abstract}
We report results on Toads and Frogs, most of them machine-checked in the Lean~4 proof
assistant. (1)~Thanatipanonda's Conjecture~T1 holds: we determine the value of every position
$\E^l\Tt^a\E^b\Ff\E^r$ with one frog and $l\le r$, as a special case of a theorem on one-frog
positions with the toads anywhere; Erickson had proved the cases $l=r=0$ and $b=0$. The other results concern the starting positions
$\Tt^a\E^b\Ff^a$. (2)~Thanatipanonda's Conjecture~T2 holds:
$\Tt\Tt\E^b\Ff\Ff=\ast$ if $2\le b\le 4$ or $b\equiv 1\pmod 6$, and $0$ otherwise, for every $b$.
(3)~$\Tt^a\E^4\Ff^a$ is infinitesimal for every $a$, and equals $\ast$ for every $a\ge5$; this is
the case $b=4$ of Conjecture~T3 and of Erickson's Conjecture~7.5. (4)~Conjecture~T4 is false:
$\Tt^6\E^{12}\Ff^6$ is neither $0$ nor $\ast$ (proved), and neither is $\Tt^7\E^{10}\Ff^7$
(computed). Both are still infinitesimal, as Erickson's conjecture predicts. We also give an
independent machine-checked proof that $\Tt^a\E^3\Ff^a=\ast$ for $a\ge4$ (the diagonal case of
Theorem~5.6.1 of~\cite{Th08}), the values of 37 further starting positions beyond the published
tables, and an account of why the method that settles $a=2$ does not reach $a=3$.
\end{abstract}

\section{Summary}

We use the numbering of the journal version of~\cite{T11}: Conjectures T1--T5 there are
Conjectures 1--5 in the arXiv preprint and in the thesis~\cite{Th08}. Erickson's conjectures are
numbered 7.1--7.6 as in~\cite{E96}.

\begin{center}
\renewcommand{\arraystretch}{1.25}
\begin{tabularx}{\linewidth}{@{}>{\raggedright}p{5.9cm}X@{}}
\toprule
Result & Status \\
\midrule
Conjecture T1: values of the one-frog positions $\E^l\Tt^a\E^b\Ff\E^r$, $l\le r$
  (Theorems~\ref{thm:onefrog}, \ref{thm:T1})
  & Proved in Lean; kernel only. \\
Conjecture T2: exact values of $\Tt\Tt\E^b\Ff\Ff$ for every $b$ (Theorem~\ref{thm:T2})
  & Proved in Lean; kernel only. \\
$\Tt\Tt\E^b\Ff\Ff$ infinitesimal for every $b$ (Erickson 7.5, $a=2$)
  & Proved in Lean; kernel only. Also a corollary of T2. \\
$\Tt^a\E^3\Ff^a$ infinitesimal for every $a$; $=\ast$ for $a\ge 4$
  & Known~\cite[Thm.~5.6.1]{Th08}; independent proof, kernel only. \\
$\Tt^a\E^4\Ff^a$ infinitesimal for every $a$ (Erickson 7.5, $b=4$)
  & Proved in Lean; the finite check uses \lean{native\_decide}. \\
$\Tt^a\E^4\Ff^a=\ast$ for every $a\ge5$ (Conjecture T3, $b=4$)
  & Proved in Lean; the finite check uses \lean{native\_decide}. \\
$\Tt^6\E^{12}\Ff^6\notin\{0,\ast\}$, refuting Conjecture T4
  & Proved in Lean; the finite check uses \lean{native\_decide}. \\
$\Tt^7\E^{10}\Ff^7\notin\{0,\ast\}$
  & Computed by exhaustive search. \\
37 new values of $\Tt^a\E^b\Ff^a$ (Table~\ref{tab:values})
  & Computed by exhaustive search. \\
\bottomrule
\end{tabularx}
\end{center}

\noindent
Every computed value agrees with all 114 published values within reach of our solver, including
those of~\cite{T11}. The four cells that Figure~2 of the preprint of~\cite{T11} lists only as
``$\neq 0$'', namely $(a,b)=(5,11),(6,8),(7,6),(7,7)$, are all~$\ast$.

\paragraph{On trust.}
A Lean proof is checked by Lean's small trusted kernel. Where a theorem needs a very large finite
computation, we evaluate it with \lean{native\_decide}, which compiles the computation to machine
code and runs it; such a theorem trusts Lean's compiler as well as the kernel, and \lean{\#print
axioms} lists one extra axiom per use. We say which results are of which kind. The main
kernel-only statements (Theorems~\ref{thm:T1} and~\ref{thm:T2}, Erickson's exceptional case
$\Tt\Tt\Tt\E\E\Ff\Ff\Ff=\pm\frac18$, and the reduction of Erickson's Conjecture~7.5 in
Section~\ref{sec:erickson}) have also been checked with Comparator, a tool which confirms that a
proof establishes exactly the statements written in a short, self-contained file
(\texttt{Challenge.lean}) using only the standard axioms. To trust these results one needs to read
only that file, which defines games, Conway's order, infinitesimals, $\ast$, the rules of the game
and the positions concerned from scratch.

\section{Conventions}

Left moves toads $\Tt$ rightwards and Right moves frogs $\Ff$ leftwards; a piece slides into an
adjacent empty square or jumps one opposing piece into the empty square beyond. We write
$G\fuzzy H$ for confused (incomparable) games.

The start $G=\Tt^a\E^b\Ff^a$ is its own mirror image with colours swapped, so $G=-G$. Three easy
consequences are used throughout:
\begin{itemize}
\item $G\le 0$ implies $G=0$, and $G\le\ast$ implies $G=\ast$ (since $-\ast=\ast$);
\item $G$ is infinitesimal if and only if its Left stop is $\le 0$ (Right stop $\ge 0$ is then
  automatic). In the formalisation, ``Left stop $\le 0$'' is replaced by an equivalent recursion
  with no numbers in it: Right may answer a Left move by \emph{passing}, once, provided the
  position is then $\le 0$ in normal play.
\item Every proof in Sections~\ref{sec:T2}--\ref{sec:T4} is a strategy for Right showing that some position is $\le 0$ (or
  $\le\ast$, or has Left stop $\le0$). The computer finds the strategy; Lean checks that it wins
  on every board size.
\end{itemize}

\section{Conjecture T1: positions with one frog}\label{sec:T1}

Conjecture~T1 of~\cite{T11} gives the values of the positions $\E^l\Tt^a\E^b\Ff\E^r$ with
$l\le r$. Two families of these were proved by Erickson~\cite{E96}: the starting positions
$\Tt^a\E^b\Ff$ ($l=r=0$; his Theorem~5.2), and the blocks $\E^l\Tt^a\Ff\E^r$ ($b=0$, every $l$
and $r$; his Lemmas~4.3 and~4.4). The conjecture joins the two: a gap $b\ge1$ together with empty
squares at the ends. We prove it as a special case of a statement about every position with one
frog, with the toads anywhere.

\subsection{The general theorem}

Consider a position with exactly one frog and at least one toad to its left; the toads to its
right have been \emph{passed}. Number the squares $0,\dots,N-1$; let the toads be at
$x_1<\dots<x_k$, the frog at $f$, and let $m$ toads lie left of the frog. Define
\begin{itemize}
\item $\varphi=\sum_{i=1}^{k}(N-k+i-1-x_i)-m-f$: the number of Left's moves minus the number of
  Right's moves if every toad left of the frog jumps it and the frog walks to the left end;
\item the \emph{room} $\rho$: the number of empty squares right of the frog minus the number of
  empty squares left of the toad nearest the frog on its left;
\item the \emph{gap} $d$: the number of empty squares between that toad and the frog.
\end{itemize}

\begin{theorem}[Lean: \lean{value\_room\_*}]
\label{thm:onefrog}
The value $G$ of such a position is
\[
G=\begin{cases}
\varphi & \text{if } \rho\ge2,\\
\varphi & \text{if } \rho=1 \text{ and } d \text{ is even},\\
\varphi+\tfrac12 & \text{if } \rho=1 \text{ and } d \text{ is odd},\\
\varphi+1 & \text{if } \rho=0,\ d \text{ is even and } \varphi\ge-1,\\
(\varphi+1)\ast=\{\varphi+1\mid\varphi+1\} & \text{if } \rho=0,\ d \text{ is odd and } \varphi\ge0.
\end{cases}
\]
\end{theorem}

{\raggedright In Lean the five cases are \lean{value\_room\_ge\_two}, \lean{value\_room\_one\_even},
\lean{value\_room\_one\_odd}, \lean{value\_room\_zero\_even} and
\lean{value\_room\_zero\_odd}.\par}

For $\rho<0$ the values are not this regular (with $\rho=-1$ and $d$ odd they include
$\varphi+\frac32$, $\varphi+\frac74$ and $\{\varphi+2\mid\varphi+2\}$); the proof needs only bounds
there. For $\E^l\Tt^a\E^b\Ff\E^r$ we have $\rho=r-l$, $d=b$ and
$\varphi=(a-1)(b-1+r)+(r-l-1)$, which gives:

\begin{theorem}[Conjecture T1; Lean: \lean{conjectureT1}]\label{thm:T1}
Let $a\ge1$, $l\le r$ and $n=(a-1)(b-1+r)$. Then
\[
\E^l\Tt^a\E^b\Ff\E^r=
\begin{cases}
(r-l-1)+n & \text{if } r-l\ge2,\\
n & \text{if } r-l\le1 \text{ and } b \text{ is even},\\
n+\tfrac12 & \text{if } r-l=1 \text{ and } b \text{ is odd},\\
n\ast & \text{if } r=l \text{ and } b \text{ is odd},
\end{cases}
\]
with two exceptions: $\Tt^a\E\Ff=\{\{a-2\mid1\}\mid0\}$ for $a\ge2$, and $\Tt^a\Ff=0$.
\end{theorem}

The first exception is part of the conjecture, and is in~\cite{E96} (from~\cite{WW}). The second
is not: for $l=r=b=0$ and $a\ge2$ the
formula $(a-1)(b-1+r)$ is negative, and the conjecture, stated in the preprint of~\cite{T11} for
$b\ge0$, fails there as printed. (We read the statement in the preprint; we were not able to
compare the wording of the journal version.) The proofs are checked by the kernel, with no
\lean{native\_decide}, and Theorem~\ref{thm:T1} has been checked with Comparator.

\subsection{Proof}

Three facts make the problem tractable.
\begin{itemize}
\item \emph{The frog never has a choice.} It can slide if the square on its left is empty, or
  jump if that square holds a toad and the next one is empty. Right only decides whether to move.
\item \emph{Comparing with a number is a one-step recursion.} For an integer $n\ge1$, $G\le n$
  if and only if every Left option $G^L$ satisfies $G^L\le n-1$ or has a Right option
  $G^{LR}\le n$; and $G\le 0$ if and only if every $G^L$ has a Right option $G^{LR}\le0$. There
  are similar rules for $G\ge n$ and for $n+\frac12$ and $n\ast$. No sums of games are needed.
\item \emph{Every move changes $\varphi$ by a fixed amount.} A Left move lowers $\varphi$ by $1$,
  a slide of the frog raises it by $1$, and a jump of the frog raises it by $3$: the frog has
  spent one move on two squares, and the toad it jumped will not jump. Jumping is bad for the
  player who jumps.
\end{itemize}

\paragraph{Lower bounds.} (i)~$G\ge\varphi$ for every one-frog position with $\varphi\ge0$. If
Left has no move, the position is $\E^j\Tt^i\Ff\Tt^{i'}$, and $\varphi=-(i+j)\le0$. So Left has a
move whenever $\varphi>0$, and any move keeps the bound. (ii)~For $k=1,2$: $G\ge\varphi+k$ on the
set $P_k$ of positions with $\rho\le-k$, or $\rho=1-k$ and $d$ even. Left slides the nearest toad
when $d\ge1$, which lowers $\rho$ by $1$; when $d=0$ he makes any other slide, or jumps when no
toad can slide. After a slide of the frog he returns to $P_k$ in the same way; after a jump of
the frog, $\varphi$ has risen by $3$ and (i) suffices. (iii)~The bounds $\varphi+\frac12$
($\rho=1$, $d$ odd) and $(\varphi+1)\ast$ ($\rho=0$, $d$ odd) follow in one step: the frog's slide
leads to a position covered by~(i), and Left's slide of the nearest toad to one in $P_1$ or~$P_2$.

\paragraph{Upper bounds.} These are proved together, by induction on the position:
\begin{center}
\begin{tabular}{@{}ll@{}}
\toprule
Positions & Bound \\
\midrule
no toad left of the frog & $\max(\varphi,0)$ \\
$\rho\ge2$, or $\rho=1$ and $d$ even & $\varphi$ \\
$\rho=1$, $d$ odd & $\varphi+\frac12$ \\
$\rho=0$, $d$ even & $\max(\varphi+1,0)$ \\
$\rho=0$, $d$ odd, $\varphi\ge0$ & $(\varphi+1)\ast$ \\
$\rho=-1$, $d=0$, $\varphi\ge-1$ & $\varphi+2$ \\
\bottomrule
\end{tabular}
\end{center}
and two degenerate families with $\varphi<0$. Right answers a Left move with a slide of the frog,
or not at all; he never jumps. A slide of the nearest toad lowers $\rho$ by $1$ and a slide of the
frog raises it by $1$; a jump by the nearest toad followed by a slide of the frog raises $\rho$ by
the distance to the next toad; any other Left move leaves $\rho$ and $d$ alone and simply costs
Left a move. For example, if $\rho=1$ and $d$ is even, a slide of the nearest toad gives $\rho=0$
with $d$ odd, and the frog's slide restores $\rho=1$ with $d$ even and the same~$\varphi$. The
cases with $\varphi<0$ use the inequality $\varphi\ge\rho-1$, valid when there is an empty square
right of the frog.

\paragraph{The exception.} In $\Tt^a\E\Ff$ ($a\ge2$), Left's only move is to
$\Tt^{a-1}\E\Tt\Ff=\{\Tt^{a-2}\E\Tt\Tt\Ff\mid\Tt^{a-1}\Ff\Tt\E\}=\{a-2\mid1\}$: in the first option
the frog can never move, and the second is $1$ by Theorem~\ref{thm:onefrog}. By the same theorem, Right's option
$\Tt^a\Ff\E$ is $0$.

\paragraph{How it was found.} The two integer recursions above compute, for every one-frog
position, the least integer above and the greatest integer below its value without canonical
forms. Compared with $\varphi$ on all positions of at most 14 squares, the lower bound $\varphi$
never fails and the upper bound depends only on $\rho$ once $\rho\ge2$. A program then checked
every step of the induction, as the proof uses it, on all positions of at most 15 squares
(1.09 million claims) before the formalisation; its first run failed only at positions with
$\varphi=-1$, which is how the degenerate families were found.

\section{Conjecture T2}\label{sec:T2}

\begin{theorem}[Lean: \lean{exactValuesA2}]\label{thm:T2}
For every $b\ge0$,
\[
\Tt\Tt\E^b\Ff\Ff =
\begin{cases}
\ast & \text{if } 2\le b\le 4 \text{ or } b\equiv 1 \pmod 6,\\
0 & \text{otherwise.}
\end{cases}
\]
In particular Conjecture~T2 of~\cite{T11} holds, and $\Tt\Tt\E^b\Ff\Ff$ is infinitesimal for
every~$b$.
\end{theorem}

The proof is checked by the kernel, with no \lean{native\_decide}; its axioms are the standard
\lean{propext}, \lean{Quot.sound} and \lean{Classical.choice}.

\subsection{Reduction to four-piece claims}

Let $G=\Tt\Tt\E^b\Ff\Ff$ with $b\ge 2$. Each player has a single move: Left to
$\Tt\E\Tt\E^{b-1}\Ff\Ff$ and Right to $\Tt\Tt\E^{b-1}\Ff\E\Ff$. Hence
\begin{itemize}
\item $G=0$ if and only if Right's only reply to Left's move, $\Tt\E\Tt\E^{b-2}\Ff\E\Ff$, is
  $\le0$;
\item $G=\ast$ if $\Tt\Tt\E^{b-1}\Ff\E\Ff\le 0$: Left's option is $\le 0$ (shown below), so $G$
  is not $\ge\ast$, and Right's option is $\le0$, so $G$ is not $\ge 0$; together with $G=-G$
  this gives $G=\ast$.
\end{itemize}
(That Left's option $\Tt\E\Tt\E^{m}\Ff\Ff$ is $\le 0$ for every $m\ge1$ is proved by the shape
method of Section~\ref{sec:shapes}, over 1,543 shapes; it also gives directly that $G$ is
infinitesimal.) So both values come down to claims ``$X\le0$'' about positions with two toads and
two frogs. Such a position is a piece order together with five gaps,
\[
\E^{g_0}\,c_1\,\E^{g_1}\,c_2\,\E^{g_2}\,c_3\,\E^{g_3}\,c_4\,\E^{g_4}.
\]

\subsection{Shapes and a margin lemma}\label{sec:shapes}

Pair each gap with its mirror image ($g_0$ with $g_4$, $g_1$ with $g_3$). The \emph{shape} of a
four-piece position records the piece order, the two mirrored differences $g_4-g_0$ and
$g_3-g_1$ exactly, and $\min(g_0,g_4)$, $\min(g_1,g_3)$ and $g_2$ capped at $4$: below $4$ they
are kept, above it only their parity.

The key lemma is a \emph{margin} argument. Say that $p$ and $q$ have the same shape with margin
$j$ if they agree in piece order and mirrored differences, and each capped quantity is either
equal in both or at least $4+j$ in both with the same parity. A move changes each gap by at most
one, and legality depends only on which gaps are zero; so if $p,q$ agree with margin $j+1$, every
move of a given piece in $p$ is available in $q$, and the results agree with margin~$j$. Each
shape has at most eight \emph{witnesses} (normal forms) with margin~$2$ to all its positions.
Consequently:

\begin{lemma}[Lean: \lean{le\_zero\_of\_good}]
Let $S$ be a finite set of shapes such that, for every witness $w$ of every shape in $S$ and
every Left move $w\to q$, some Right move $q\to r$ leads to a shape in $S$. Then every
position whose shape lies in $S$ is $\le0$.
\end{lemma}

The hypothesis is a finite computation, which the Lean kernel performs itself; the program that
proposes $S$ is not trusted.

\paragraph{Even $b$.} For even $k\ge4$, Right wins $\Tt\E\Tt\E^k\Ff\E\Ff$ while keeping both
mirrored differences within $4$; the positions reached fall into 1,038 shapes, found by a search
on shapes with backtracking. This gives $\Tt\Tt\E^b\Ff\Ff=0$ for even $b\ge6$. The cases $b\le30$
are checked one by one by the kernel.

\paragraph{Odd $b$: why shapes fail.} For odd $b$ the smallest bound on the mirrored differences
with which Right still wins grows with the board. For $\Tt\E\Tt\E^k\Ff\E\Ff$ with
$k=3,7,9,13,15,19$ it is $4,6,8,10,12,14$; for $\Tt\Tt\E^k\Ff\E\Ff$ with $k=6,12,18,24,30$ it is
$5,9,13,17,21$. So no finite set of shapes can work, and odd $b$ needed a different invariant.

\subsection{Odd \texorpdfstring{$b$}{b}: a linear invariant}

\paragraph{The approach phase, and where the period 6 comes from.} While the pieces are in the
order $\Tt\Tt\Ff\Ff$, Right always moves his \emph{front} frog. Each Left move lowers the race
count (empty squares in front of the toads minus those in front of the frogs) by one, and
Right's reply restores it. The rear frog never moves, so the positions stay in the
two-parameter family
\[
\E^{g_0}\,\Tt\,\E^{g_1}\,\Tt\,\E^{g_2}\,\Ff\,\E^{2g_0+g_1+c}\,\Ff,
\qquad c\in\{0,1\},
\]
with $c=0$ starting from $\Tt\E\Tt\E^{k}\Ff\E\Ff$ (the claim for $G=0$) and $c=1$ from
$\Tt\Tt\E^{k}\Ff\E\Ff$ (the claim for $G=\ast$). Exact play shows that every position of the family is a loss for
Left moving first when the number $N$ of empty squares satisfies $N\equiv 3,5\pmod 6$ (for
$c=0$) or $N\equiv1\pmod6$ (for $c=1$). Since $N=b$ in both cases, this is precisely the pattern
of Conjecture~T2 for odd $b$. (In the proof this is not used directly; it is a consequence of the
invariant below.)

\paragraph{After the first contact.} The positions after the first jump are cold: almost all
their canonical values are numbers or numbers plus~$\ast$. Accordingly, for each order of the
pieces, the set of positions that Left (moving first) loses is described by a small decision
tree over linear forms in the gaps, residues mod~$2$ and~$3$, and the race count. The trees,
learned from exact play on all positions with up to 30--40 empty squares, have
$2$ ($\Ff\Ff\Tt\Tt$), $6$ ($\Ff\Tt\Ff\Tt$), $16$ ($\Ff\Tt\Tt\Ff$), $32$ ($\Tt\Ff\Ff\Tt$) and $49$
($\Tt\Ff\Tt\Ff$) leaves, and $19$ for the contact region of $\Tt\Tt\Ff\Ff$. Right always moves
into these sets.

\paragraph{The proof.} The union of these sets is an invariant \lean{inv2} consisting of 67 cases,
each a conjunction of linear inequalities, equalities and congruences in the five gaps. The
starting positions lie in it, and it is \emph{closed}: from every position in it, every Left move
has a Right reply back into it, for every board size. Closure was first checked with the SMT
solver Z3, with the gaps as unbounded integers; Z3 also split each (piece order, Left move, case)
into regions with a single Right reply and a single target case, 379 regions in all. In Lean,
each region is a lemma in Presburger arithmetic proved by the decision procedure \lean{omega}
(a few needed a case split on a gap or its parity). Closure then gives $\le 0$ by induction on
the number of moves remaining.

\section{The columns \texorpdfstring{$b=3$}{b=3} and \texorpdfstring{$b=4$}{b=4}}

\subsection{Runs instead of gaps}

With few empty squares, the roles of pieces and gaps swap: a position is a few empty squares
separated by long \emph{runs} of pieces, and the number of runs is not fixed. Positions are
handled as run lists, and moves on run lists are proved equal to moves on cells. Erickson's dead
pieces are essential: a run enclosed by a toad run of length $\ge2$ before it and a frog run of
length $\ge2$ after it, with no empty square between, never moves again, so its length is
dropped. The \emph{shape} of a position records each run's kind, the lengths of empty runs and of
runs of length one, the sorted order of the remaining live runs of length $\ge2$ (together with an
anchor of value $2$), and the gaps between consecutive sorted lengths, exactly up to a threshold
and coarsely above it. A move changes each recorded difference by at most~$2$, which gives a margin
lemma as in Section~\ref{sec:shapes}. The framework is proved once in Lean for any such gap rule.

\subsection{\texorpdfstring{$b=3$}{b=3}}

\begin{theorem}[Lean: \lean{infinitesimal\_start\_b\_three}, \lean{start\_three\_value}]
$\Tt^a\E^3\Ff^a$ is infinitesimal for every $a$, and equals $\ast$ for every $a\ge4$ (and $0$ for
$a=0,3$, $\ast$ for $a=1,2$).
\end{theorem}

This is known: it is the diagonal case of Theorem~5.6.1 of~\cite{Th08},
$\Tt^a\E^3\Ff^b=\{a-b\mid a-b\}$ for $a,b\ge4$ (Theorem~D.3 of the preprint of~\cite{T11}). Our proof
is independent and, as far as we know, the first machine-checked one. Right's strategy uses 3,031
shapes for infinitesimality and 1,263 for the value; gaps are kept exactly up to~$2$ and by parity
above. The finite checks are done by the kernel. We have not formalised the off-diagonal case
$a\neq b$.

\subsection{\texorpdfstring{$b=4$}{b=4}}

\begin{theorem}[Lean: \lean{infinitesimal\_start\_b\_four}, \lean{start\_four\_value};
\lean{native\_decide}]\label{thm:b4}
$\Tt^a\E^4\Ff^a$ is infinitesimal for every $a$. Moreover $\Tt^a\E^4\Ff^a=\ast$ for every
$a\ge5$, with $0$ for $a=0,1,4$ and $\ast$ for $a=2,3$.
\end{theorem}

This is the case $b=4$ of Conjecture~T3 and of Erickson's Conjecture~7.5. Values were known for
$a\le10$~\cite{T11}, and~\cite[p.~89]{Th08} lists $\Tt^a\E^4\Ff^b$ as open.

The settings that work for $b=3$ fail here: Right needs up to $13$ runs, and with parity kept the
search for a closed set of shapes does not terminate, wandering through positions with four or
five long runs that differ only in the order, spacing and parity of their lengths. The margin
argument never uses parity, and dropping it closes the search. For infinitesimality, gaps between
long runs are kept exactly up to~$1$ and otherwise only as ``far''; Right's strategy uses 233,087
shapes and 2.8 million witnesses. For the value, gaps are kept exactly up to~$2$; Right's strategy
uses 361,981 shapes and 3.5 million witnesses. These checks are about 270 times larger than those
of $b=3$; the kernel would need an estimated 150 CPU hours, so they are evaluated with
\lean{native\_decide}. As independent checks, two separate programs confirm that the shape sets
are closed, and Right's strategy has been played on the actual game trees of $\Tt^a\E^4\Ff^a$ for
every $a\le24$ (for the value, $5\le a\le22$); each of these is a complete proof for that~$a$,
independent of the margin argument.

\section{Counterexamples to Conjecture T4}\label{sec:T4}

Conjecture~T4 of~\cite{T11} states that $\Tt^a\E^b\Ff^a$ is $0$ or $\ast$ whenever $b\ge a>0$.
The table in~\cite{T11} stops at $b=8$ for $a=6$ and at $b=7$ for $a=7$.

\begin{theorem}[Lean: \lean{not\_thanatipanonda\_conjecture4}; \lean{native\_decide}]
$G=\Tt^6\E^{12}\Ff^6$ satisfies $G\fuzzy 0$ and $G\fuzzy\ast$. In particular $G$ is neither $0$
nor~$\ast$, and Conjecture~T4 is false.
\end{theorem}

\begin{proof}[Outline]
Right has a single move from $G$, the front frog's slide, to
$R=\Tt^6\E^{11}\Ff\E\Ff^5$. We show $R\le0$ and $R\le\ast$ by a strategy table for Right of about
21 million positions, extracted from an exact solver and checked in Lean (with
\lean{native\_decide}). Since $G$ has a Right option $R\le 0$, $G\not\ge0$; since $G+\ast$ has the
Right option $R+\ast\le0$, $G\not\ge\ast$. As $G=-G$, also $G\not\le0$ and $G\not\le\ast$.
\end{proof}

By exhaustive search (not formalised), $\Tt^7\E^{10}\Ff^7$ is likewise neither $0$ nor $\ast$.
Both positions are nonetheless infinitesimal (their Left stops are $0$), in agreement with
Erickson's Conjecture~7.5, and neither is $\ast2$ or $\ast3$. We have not determined their
canonical forms. Outside the column $b=2$, they are the only cells of Table~\ref{tab:values} whose value is
neither $0$ nor~$\ast$. Not to be confused with Erickson's
Conjecture~7.4, which~\cite{T11} refutes.

\section{Values of the starting positions}\label{sec:values}

Table~\ref{tab:values} collects every known value of $\Tt^a\E^b\Ff^a$ for $a\le14$, $b\le30$.
Our values come from an exhaustive solver (\texttt{tools/values.cpp}) on boards of up to 32
cells. Like the program of~\cite{T11}, it decides comparisons rather than computing canonical
forms: since $G=-G$, deciding $G\le0$ and $G+\ast\le0$ determines whether $G$ is $0$, $\ast$ or
neither.

\begin{table}[htbp]
\centering
\small
\setlength{\tabcolsep}{4.2pt}
\renewcommand{\arraystretch}{1.08}
\begin{tabular}{@{}r|*{14}{c}@{}}
\toprule
$b\backslash a$ & 1 & 2 & 3 & 4 & 5 & 6 & 7 & 8 & 9 & 10 & 11 & 12 & 13 & 14 \\
\midrule
\csname @@input\endcsname values_table.tex
\bottomrule
\end{tabular}
\caption{Values of $\Tt^a\E^b\Ff^a$ (rows $b$, columns $a$; note the transposition relative to
Figure~2 of~\cite{T11}). \colorbox{orange!22}{Shaded}: new, computed here (including the four
``$\neq0$'' cells of~\cite{T11}; the rows $b\le3$ for $a\ge11$ are not counted, since they
follow from~\cite{T11,Th08}). \textbf{Bold}: proved in Lean (the column $a=2$, the rows $b=3,4$;
for $(6,12)$, that the value is neither $0$ nor $\ast$). $\mathrm{N}$: an infinitesimal with a long canonical form, as in~\cite{E96,T11}.
$\dagger$: infinitesimal, but neither $0$, $\ast$, $\ast2$ nor $\ast3$. Blank: not computed.
Beyond the table, $a=2$ continues with period~$6$ (Theorem~\ref{thm:T2}), and~\cite{T11} gives
$a=3$, $31\le b\le 53$: all~$0$.}
\label{tab:values}
\end{table}

Some observations on the new values:
\begin{itemize}
\item In the region $b\ge a$ of Conjecture~T4, every new value other than the two
  counterexamples is $0$ or $\ast$. The value $\ast$ keeps recurring for $a=5$ ($b=17,21$) and
  $a=6$ ($b=15$), so for $c=5,6$ the data are too short to suggest the threshold $b_0$ of
  Conjecture~T5.
\item For $a=4$ the values are $0$ for every $b$ from $10$ to $24$.
\item In the region $a>b>0$, $b\ne2$, of Conjecture~T3, every computed value is $\ast$.
\end{itemize}

\section{Erickson's Conjecture 7.5}\label{sec:erickson}

Erickson's Conjecture~7.5 remains open. Combining the families above with finitely many checked
cases, we have in Lean:
\begin{itemize}
\item infinitesimality for $a\le2$ (every $b$), $b\le3$ (every $a$, with $b=2$ due
  to~\cite{T11} and formalised here), and every case with $2a+b\le12$; Erickson's exception
  $\Tt\Tt\Tt\E\E\Ff\Ff\Ff=\pm\frac18$ is also proved;
\item hence the conjecture is equivalent to its restriction to $a\ge3$, $b\ge4$, $2a+b\ge13$
  (\lean{conjecture5\_of\_openCases}, kernel only, checked with Comparator);
\item in addition, $b=4$ for every $a$ (Theorem~\ref{thm:b4}, \lean{native\_decide}), and the cases
  $(3,7)$ to $(3,12)$ and $(4,5)$ by strategy certificates.
\end{itemize}
By exhaustive search, the conjecture holds on every board of at most $25$ squares, i.e.\ for all
$2a+b\le25$.

\section{The row \texorpdfstring{$a=3$}{a=3}: a negative result}\label{sec:a3}

The natural next step, and the route towards Conjecture~T5 for $c=3$, is the row $a=3$. As for
$a=2$, it reduces to a claim about Left's unique first move: $\Tt\Tt\E\Tt\E^m\Ff\Ff\Ff\le0$ for
$m\ge5$. We report what we found, since it may save effort.

\begin{itemize}
\item \emph{The approach phase works.} If Right always moves his front frog while the centre gap
  is at least~$2$, the race count stays at $-1$ and the positions form the three-parameter family
  $\E^{g_0}\Tt\E^{g_1}\Tt\E^{g_2}\Tt\E^{g_3}\Ff\E^{3g_0+2g_1+g_2-1}\Ff\Ff$. Every position of this
  family is a loss for Left moving first for every number of empty squares from $6$ to $34$,
  with no congruence condition (all are wins at $5$, matching the failure of the claim at $m=4$).
\item \emph{After contact it does not.} All 19 later piece orders occur, and the exact loss sets
  have no small description: decision trees learned on boards of up to 18 empty squares have
  hundreds to over a thousand leaves and fail on larger boards. A search for a smaller closed
  subset of the losses (growing the set whenever a Left move has no answer in it) always closes
  on the data, but needs about $930$ linear cases for boards of at most $14$ empty squares,
  $1{,}760$ for at most $18$, and $3{,}589$ for $22$ to $30$; it does not level off, and fails just
  outside the sizes on which it was found. For $a=2$ the whole invariant had 67 cases.
\item \emph{No spare tempo.} $\Tt\Tt\E\Tt\E^m\Ff\Ff\Ff$ is a number:
  \[
  \begin{array}{c|cccccccccccccc}
  m & 5 & 6 & 7 & 8 & 9 & 10 & 11 & 12 & 13 & 14 & 15 & 16 & 17 & 18\\\hline
  \text{value} & 0 & -\frac18 & -\frac18 & -\frac18 & -\frac14 & -\frac14 & -\frac14 & -\frac14
    & 0 & -\frac14 & -\frac12 & -\frac14 & -\frac14 & -\frac12
  \end{array}
  \]
  So Left's first move costs him only a fraction of a move, and not uniformly: at $m=13$, where
  $\Tt^3\E^{14}\Ff^3=\ast$, the value is exactly~$0$.
\item \emph{Structure.} Right must jump: restricted to at most four jumps of his own, he
  loses for most $m\le22$. When Left refuses to jump, Right wins by blocking toads rather than
  through the race count. On boards of 45--50 empty squares, two-dimensional
  sections of the loss sets look regular: polygons with edges of slope $1$, $2$ and $\frac12$, and
  thin stripes of period~$6$.
\end{itemize}

\noindent
Our reading is that a proof of the row $a=3$ needs a strategy for Right chosen for its regularity
after contact as well, or an invariant language in which the period-6 stripes and the blocking of
toads are primitive. We would be glad to hear of any structural insight into this row.

\section{Status of the conjectures}

\begin{center}
\renewcommand{\arraystretch}{1.25}
\begin{tabularx}{\linewidth}{@{}l>{\raggedright\arraybackslash}X@{}}
\toprule
Conjecture & Status \\
\midrule
Erickson 7.1, 7.3 & True: proved by Thanatipanonda (Theorems D.2, D.1 of the preprint
  of~\cite{T11}; proofs in~\cite{Th08}). \\
Erickson 7.2 & True: proved in~\cite{TZ07}. \\
Erickson 7.4 & False: $\Tt^7\E^7\Ff^6>2$~\cite{T11}. \\
Erickson 7.5 & Open. New: $a=2$ and $b=4$ for every $a$ (Sections~\ref{sec:T2},
  \ref{sec:erickson}). \\
Erickson 7.6 & True (Hull, 2000, as reported in~\cite{T11}). \\
Thanatipanonda T1 & \textbf{True} (Theorem~\ref{thm:T1}; the cases $l=r=0$ and $b=0$ are
  in~\cite{E96}), with $\Tt^a\Ff=0$ in place of the
  formula when $l=r=b=0$. \\
Thanatipanonda T2 & \textbf{True} (Theorem~\ref{thm:T2}). \\
Thanatipanonda T3 & Open. Known for $b=1$~\cite{WW} and $b=3$~\cite{Th08};
  \textbf{new: $b=4$} (Theorem~\ref{thm:b4}). All computed values agree. \\
Thanatipanonda T4 & \textbf{False}: $(6,12)$ (proved) and $(7,10)$ (computed). \\
Thanatipanonda T5 & Open. For $c=3$ the last $\ast$ is at $b=29$~\cite{T11}; the row $a=3$
  resists the methods above (Section~\ref{sec:a3}). \\
\bottomrule
\end{tabularx}
\end{center}

\paragraph{Literature.} We call a result new when we could not find it in~\cite{E96,TZ07,T11,Th08},
in the ``Unsolved problems in combinatorial games'' lists of the six \emph{Games of No Chance}
volumes, in the works citing these (via OpenAlex, Crossref, Semantic Scholar and zbMATH Open), on
arXiv, or in formal libraries (Mathlib and other Lean, Rocq and Isabelle code). We could not reach
the preprints \emph{Three (More) Results of Combinatorial Game Toads and Frogs} and \emph{Library
values for Toads and Frogs} cited in~\cite{T11}, nor Siegel's \emph{Combinatorial Game Theory} or
\emph{Lessons in Play}; corrections are very welcome.

\paragraph{Reproducing.} Everything is in the repository named on the title page. The default
build (\texttt{lake build}, Lean \texttt{v4.36.0-rc2}) contains the kernel-only results; the
\lean{native\_decide} results are built by \texttt{lake build CgtNative} (about 15 minutes for
$b=4$); the comparator check is \texttt{lake comparator}. The solvers and search programs are in
\texttt{tools/} and \texttt{scripts/}, and fuller accounts are in \texttt{docs/ONE\_FROG.md},
\texttt{docs/ROW\_A2.md},
\texttt{docs/COLUMN\_B3.md} and \texttt{docs/RESEARCH\_NOTES.md}.

\begin{thebibliography}{TZ07}
\bibitem[E96]{E96} J.~Erickson, \emph{New Toads and Frogs results}. In R.~J.~Nowakowski (ed.),
  \emph{Games of No Chance}, MSRI Publications~29, Cambridge University Press, 1996, 299--310.
\bibitem[T11]{T11} T.~Thanatipanonda, \emph{Further hopping with Toads and Frogs}. Electron.\ J.\
  Combin.\ 18(1) (2011), \#P67. Preprint arXiv:0804.0640 (2008).
\bibitem[Th08]{Th08} T.~Thanatipanonda, \emph{Symbolic-computational methods in combinatorial game
  theory and Ramsey theory}. PhD thesis, Rutgers University, 2008. doi:10.7282/T30G3K3F.
\bibitem[TZ07]{TZ07} T.~Thanatipanonda and D.~Zeilberger, \emph{A symbolic finite-state approach
  for automated proving of theorems in combinatorial game theory}. J.\ Difference Equ.\ Appl.\
  15(2) (2009); arXiv:0710.4951.
\bibitem[WW]{WW} E.~R.~Berlekamp, J.~H.~Conway and R.~K.~Guy, \emph{Winning Ways for your
  Mathematical Plays}. Academic Press, 1982.
\end{thebibliography}

\end{document}
